\section{Conclusion}
\label{sec:conclusion}

This work presents \textbf{ViBridge-VQA}, a parameter-efficient adaptation of the
Vietnamese VLM Vintern-1B for generative visual question answering. By keeping the
vision encoder and the language decoder frozen and training only a lightweight
multi-token bridge and a rank-16 decoder LoRA, ViBridge-VQA updates 1.01\% of the
parameters and encodes a single image view. On AutoViVQA it matches the
fine-tuned Vintern-1B on F1, achieves the best generation scores among all
compared systems, and is stable across seeds and splits. Our ablation shows that
the remaining gap is not closed by larger bridges, multi-reference training or
alignment to the original projector, whereas decoder adaptation improves every
bridge and makes them perform almost equally: for a frozen VLM of this size, the
bottleneck lies in the decoder rather than in the visual pathway.

ViBridge-VQA also has limitations. It still trails ViMoE-VQA on precision and F1,
and the fine-tuned Vintern-1B on exact-match accuracy. Because the bridge reads one
global vector of a single view, it is not suited to questions that require
reading scene text\ifoodtext, on which it falls far behind the original
Vintern-1B,\fi{} and it does not improve on counting. The baselines are reported on a different split of
AutoViVQA and were not reproduced in our setting, and our conclusions concern a
frozen 0.5B decoder and may not transfer to larger or fully trainable decoders.
Future work will combine the global-vector bridge with patch-level or
high-resolution tokens for text-heavy questions, and study larger Vietnamese
decoders under the same frozen visual pathway. The code and the grouped split are
available at \url{https://github.com/nguyennn263/modular-vlm-finetune}.

Overall, ViBridge-VQA shows that a strong Vietnamese VQA model can be obtained by
adapting a small fraction of an existing VLM, and that a careful component-wise
analysis indicates where further effort is best spent. We hope this work will
encourage efficient and well-evaluated multimodal models for low-resource
languages.
